\documentclass[conference]{IEEEtran}
\IEEEoverridecommandlockouts
\usepackage{cite}
\usepackage{amsmath,amssymb,amsfonts}
\usepackage{booktabs}
\usepackage{graphicx}
\usepackage{textcomp}
\usepackage{xcolor}
\usepackage{url}
\usepackage{array}
\def\BibTeX{{\rm B\kern-.05em{\sc i\kern-.025em b}\kern-.08em
    T\kern-.1667em\lower.7ex\hbox{E}\kern-.125emX}}

\begin{document}

\title{UNIVERSAL-CIRCULAR-RECOVERY: Quality-Aware Design and Routing for Repeated Material Service}

\author{\IEEEauthorblockN{Xcientist Research Proposal}
\IEEEauthorblockA{Survey, Idea, ExperimentDesign, and Author pipeline\\
Design-only study package}}

\maketitle

\begin{abstract}
The proposition that essentially every material could be recycled and reused is attractive but scientifically underspecified. Technical processability, product-grade recovery, actual collection, economic viability, and environmental benefit are different properties. This paper proposes UNIVERSAL-CIRCULAR-RECOVERY, a quality-aware framework that co-designs product architecture, material formulation, identity and sorting, recovery routes, markets, and policy. Reuse, repair, refurbishment, mechanical recycling, chemical recovery, biological or thermal treatment, controlled disposal, and substitution compete within the same decision system. Candidates are evaluated per delivered service rather than per tonne of waste, and a route is counted as successful only when its output meets a receiving-product specification, displaces a defined virgin resource, and remains feasible under regional contamination, infrastructure, market, and uncertainty scenarios. The design covers beverage containers, food packaging, construction products, textiles, electronics, and contaminated or hazardous products. It measures collection, sorting, quality retention, repeated cycles, residues, energy, water, emissions, toxicity, leakage, private cost, social cost, and resilience. A Pareto frontier and explicit exclusion rules replace a universal recycling score. This is a design-only proposal: no facility trial, life-cycle inventory, economic estimate, or material ranking is presented as an observed result. The direct conclusion is that near-universal circularity is plausible as a governed portfolio of designed product systems, but literal recovery of every object is neither a meaningful nor a safe engineering target.
\end{abstract}

\begin{IEEEkeywords}
material recycling, reuse, circular economy, design for recovery, sorting, contamination, life-cycle assessment, techno-economic analysis
\end{IEEEkeywords}

\section{Introduction}

Recycling technology has expanded from a small set of clean, homogeneous streams to a broad family of mechanical, chemical, thermal, biological, and robotic processes. It is now possible to recover useful fractions from many products that were previously discarded. This progress motivates a larger question: can essentially every material be recycled and reused?

The question contains a hidden change of subject. A laboratory process that transforms an object is not the same as a system that collects the object, separates it, produces a marketable output, displaces virgin material, and reduces total environmental and social burdens. A contaminated package can be technically treated but fail a food-contact specification. A composite product can be separated only with high energy and solvent use. A reusable item can require so many wash and return cycles that a single-use recyclable alternative has lower total burden. A route may also be profitable because cleanup, worker exposure, climate damage, or residue disposal is paid by someone else.

The U.S. Environmental Protection Agency defines recycling as collecting and processing materials that would otherwise be discarded into new products, and its hierarchy places source reduction and reuse before recycling \cite{epa}. The same source describes recycling as a loop involving collection and processing, manufacturing, and demand for recycled products. It identifies contamination, infrastructure, communication, product design, market development, and measurement as system challenges. These observations imply that a universal recycling machine cannot by itself create a circular economy.

This paper reformulates the topic as:

\begin{quote}
For a defined material-product stream and required service, under what combinations of design, collection, sorting, processing, market demand, and policy can a high fraction of material be repeatedly recovered and reused without increasing life-cycle environmental burdens, toxicity, or total social cost beyond acceptable bounds?
\end{quote}

The proposed answer is a portfolio architecture called UNIVERSAL-CIRCULAR-RECOVERY. It does not assume that all objects deserve the same route. It searches for material-product systems in which reduction, reuse, repair, refurbishment, mechanical recycling, chemical recovery, biological or thermal treatment, controlled treatment, or substitution retain the most future service. Its four contributions are:

\begin{enumerate}
\item a vocabulary separating technical processability, product-grade recycling, actual routing, reuse, and circularity;
\item an equal-service representation connecting object design to quality-adjusted recovery;
\item an experimental and modeling protocol for contamination, repeated processing, markets, external costs, and regional infrastructure; and
\item falsification gates that can reject attractive but incomplete recycling claims.
\end{enumerate}

The work is a design-only research proposal. It reports no completed trial or measured route ranking.

\section{Evidence and Scientific Scope}

\subsection{Material flows and the recycling hierarchy}

Historical material-flow analysis has shown that production, use, and disposal can grow faster than recovery capacity. For plastics, Geyer, Jambeck, and Law documented a large cumulative stock in landfills and the environment and a small fraction recycled into equivalent products \cite{geyer}. The exact values are boundary- and year-dependent, but the methodological lesson generalizes: circularity must be measured as a flow of products and services, not inferred from a material label.

Hopewell, Dvorak, and Kosior describe plastics recycling as a combination of technical, economic, and policy challenges \cite{hopewell}. Their framework is useful beyond plastics because the same failure sequence appears elsewhere: collection is incomplete, streams are mixed, contaminants reduce quality, processing creates residues, and recovered output competes with volatile virgin prices. A material can therefore be processable in principle while being uneconomic or environmentally inferior in practice.

Reduction, reuse, repair, and refurbishment are not competitors that should be excluded from a recycling study. Reuse retains product structure and often avoids remanufacturing, but it requires durability, return logistics, inspection, cleaning, and high enough cycle counts. Recycling can be preferable when the product is dispersed, contaminated, difficult to return, or inexpensive to process into a useful grade. The correct hierarchy is conditional on service, distance, energy, water, loss, and safety.

\subsection{Terminology and decision boundaries}

Table~\ref{tab:terms} defines the terms used in this paper. The distinctions are operational: each one has a different measurement and a different failure mode.

\begin{table}[t]
\caption{Operational material-recovery vocabulary}
\label{tab:terms}
\centering
\footnotesize
\setlength{\tabcolsep}{3pt}
\begin{tabular}{p{0.28\columnwidth}p{0.62\columnwidth}}
\toprule
Term & Operational meaning\\
\midrule
Technically recoverable & A specified process transforms the material under stated conditions.\\
Product-grade recyclable & Output meets the next product's specification and safety constraints.\\
Actually recyclable & Collected, sorted, processed, sold, and used in the modeled region.\\
Reusable & Product or component delivers another service after inspection, cleaning, repair, or refurbishment.\\
Circular & Material value is retained through repeated service cycles within a stated boundary.\\
Economically viable & Defined private or social costs are covered under a declared scenario.\\
\bottomrule
\end{tabular}
\end{table}

The primary functional unit is one delivered service. Examples include one beverage-container service, one food-protection service for a specified shelf life, one square meter-year of construction performance, one garment-year of wear, one electronic-device service, or one regulated medical or hygiene service. One tonne of waste is retained as a secondary facility unit but cannot replace equal-service comparison.

The system boundary begins at design and includes material production, additives, manufacturing scrap, distribution, use, collection, transport, sorting, cleaning, disassembly, conversion, purification, recovered-product manufacturing, residues, leakage, and avoided virgin production. Reuse adds washing, repair, return transport, loss, damage, inspection, and cycle count. Worker and community burdens are reported separately from private cash flow.

\subsection{Why every material is a boundary problem}

Materials do not remain concentrated and separable forever. Abrasion, corrosion, evaporation, combustion, biological metabolism, contamination, and dispersion can move material into dilute or chemically altered forms. Some products also contain hazardous substances or biological contamination for which open recovery is unsafe. A scientific framework should not assign a recycling target to every stream regardless of risk. Excluding a stream from open recovery is valid only when the hazard, containment route, and substitution opportunity are documented.

\section{UNIVERSAL-CIRCULAR-RECOVERY Architecture}

\subsection{Coupled product and system representation}

Each candidate system is represented as

\begin{equation}
z=(a,m,d,i,s,k,h,p,r),
\label{eq:candidate}
\end{equation}

where $a$ is product architecture, $m$ material family, $d$ additive and joining design, $i$ identity and information, $s$ sorting configuration, $k$ treatment route, $h$ reuse or recovery history, $p$ market and policy scenario, and $r$ regional context. This representation makes coatings, adhesives, pigments, flame retardants, batteries, and prior processing history explicit.

The architecture contains five linked modules. The design module chooses geometry, material families, joins, additives, labels, and access points. The identity module records composition and safe handling information in a machine-readable form. The routing module combines identity with measured condition, contamination, and local facility capability. The recovery module measures reuse, disassembly, mechanical, chemical, biological, thermal, or controlled treatment. The decision module compares service and burden vectors and reports a Pareto set.

The modules are linked through data contracts rather than one opaque score. A material property cannot be averaged with a carbon impact before service feasibility is checked. A successful reactor conversion cannot be counted as circularity unless the recovered output displaces a specified primary resource. A policy incentive changes a scenario; it does not change measured chemistry.

\subsection{Product strata}

The experimental portfolio uses contrasting streams to expose different constraints.

\begin{table*}[t]
\caption{Product strata and dominant recovery constraints}
\label{tab:strata}
\centering
\footnotesize
\setlength{\tabcolsep}{3pt}
\begin{tabular}{p{0.17\textwidth}p{0.29\textwidth}p{0.29\textwidth}p{0.19\textwidth}}
\toprule
Stratum & Example service & Candidate recovery actions & Dominant risks\\
\midrule
Beverage containers & Specified volume and number of safe uses & Deposit reuse, washing, mechanical recycling, remelting & Breakage, transport, labels, market demand\\
Food packaging & Protect a product for a specified shelf life & Mono-material recycling, fiber recovery, composting where certified & Residue, multilayers, product loss\\
Construction products & Meet a structural or thermal duty for a design life & Reuse, refurbishment, aggregate recovery, remanufacture & Demolition mixture, transport, contamination\\
Textiles and furnishings & One garment-year or furnishing-year of service & Repair, resale, fiber-to-fiber, fiber recovery & Blends, dyes, finishing chemicals, short life\\
Consumer electronics & One device service with safety and performance & Repair, component reuse, material separation, controlled battery treatment & Miniaturization, hazardous parts, data and labor\\
Medical or hazardous products & One regulated safe or sterile service & Validated reuse, segregated recovery, controlled treatment, substitution & Biohazard, toxic additives, regulatory failure\\
\bottomrule
\end{tabular}
\end{table*}

\subsection{Quality-aware routing}

Let $q_k(r)$ be the probability that an object in region $r$ enters route $k$. Routes include reuse, repair, refurbishment, mechanical recycling, chemical recovery, biological or thermal processing, controlled treatment, and disposal. Collection, sorting, processing, grade acceptance, and market acceptance determine whether route output is useful. The quality-adjusted route yield is

\begin{equation}
Y_{k,r}=q_{\mathrm{collect}}q_{\mathrm{sort}}q_{\mathrm{process}}q_{\mathrm{grade}}q_{\mathrm{market}}.
\label{eq:yield}
\end{equation}

Each factor is estimated from observed or bounded evidence. A high chemical conversion yield does not compensate for purification failure or the absence of a buyer. A mass recovery percentage without grade and market factors is reported as process yield, not product-grade circularity.

For repeated use, the delivered service multiplier is

\begin{equation}
S(c)=\frac{n_{\mathrm{cycles}}(c)[1-\lambda_{\mathrm{service}}(c)]}{1+\rho_{\mathrm{loss}}(c)},
\label{eq:service}
\end{equation}

where $n_{\mathrm{cycles}}$ is the number of successful cycles, $\lambda_{\mathrm{service}}$ is service failure probability, and $\rho_{\mathrm{loss}}$ represents replacement and loss burden. The complete inventory still includes washing, repair, transport, and inspection. This equation prevents a product with a high nominal cycle count but high failure or loss from being assigned an unrealistic benefit.

\section{Material, Design, and Facility Experiments}

\subsection{Design-for-recovery intervention}

Construct matched conventional and design-for-recovery products for each applicable stratum. Vary mono-material versus multilayer construction, reversible versus permanent joins, label and adhesive formulation, additive disclosure, component access, and machine-readable identity. Freeze the delivered service, safety, mass, and lifetime thresholds before testing. Record manufacturing scrap, assembly time, material use, disassembly time, and product performance after aging.

The intervention is not assumed to be beneficial. A simpler package may lose barrier and increase food waste. A reversible fastener may reduce durability. A digital identity layer may require sensors, electricity, and data governance. These penalties are measured and included in the same functional unit.

\subsection{Contamination and sorting}

Prepare blinded samples at controlled contamination levels: clean, moisture, food or biological residue, incompatible material, pigment or additive, label or adhesive, mixed assembly, and unknown history. Use optical, near-infrared, magnetic, density, robotic, and manual separation as appropriate. Measure false positive and false negative rates, route-specific contamination, reject composition, worker exposure controls, and the fate of material that cannot be classified.

The experiment distinguishes object-level recovery from stream-level damage. A clean item can be recovered while contaminating the output of an entire batch. Conversely, a low-value residue may be safely isolated without lowering the quality of the remaining stream. Both mass and quality are required.

\subsection{Repeated recovery and reuse}

For reusable products, conduct service cycles with realistic washing, inspection, repair, return logistics, breakage, loss, and user damage. Measure water, energy, detergents, labor, transport, and successful service. For mechanical processing, run repeated cycles and measure yield, purity, mechanical or structural property retention, odor, color, additive carryover, and next-product grade. For chemical recovery, measure conversion, depolymerization or feedstock yield, solvent and catalyst recovery, purification, energy, residue, and virgin-material displacement.

Biological and thermal routes are evaluated by receiving environment and treatment condition. Report products, emissions, residues, methane control, and containment. Fragmentation, ash production, or mass loss is not automatically counted as material recovery. A route is credited only for a defined product, feedstock, or service.

\subsection{Regional route audit}

Use at least three contrasting regions: one with high collection and established markets, one with established recycling but limited reuse or treatment capacity, and one with low collection and high leakage or informal-handling risk. For each region, measure or bound collection coverage, sorting technology, facility distance, throughput, route capacity, contamination, electricity, labor, transport, market acceptance, and residue management.

Intended labeling and actual routing are separate variables. Stress scenarios include low collection, high contamination, sorting error, facility failure, price shock, buyer withdrawal, and increased transport distance. A candidate that works only with an unavailable facility is reported as an infrastructure-dependent strategy, not as a present universal solution.

\section{Economic and Life-Cycle Formulation}

\subsection{Impact and cost vector}

For candidate $c$ and region $r$, report

\begin{equation}
J(c,r)=\left(G,E,W,T,P,C_{\mathrm{priv}},C_{\mathrm{soc}},Q,S,R\right),
\label{eq:objective}
\end{equation}

where $G$ is greenhouse-gas impact, $E$ energy, $W$ water, $T$ toxicity and exposure, $P$ persistence or leakage, $C_{\mathrm{priv}}$ private cost, $C_{\mathrm{soc}}$ social cost, $Q$ quality-adjusted circularity, $S$ delivered service, and $R$ resilience. Environmental categories may be expanded for land, biodiversity, eutrophication, air pollution, and resource depletion.

The private ledger contains collection, transport, labor, utilities, capital, maintenance, processing, product revenue, recovered-material revenue, and disposal cost. The social ledger adds external climate damage, pollution, ecological damage, cleanup, health burden, worker and community exposure, and distributional effects. The two ledgers are not averaged without an explicit policy question.

For service unit $u$, a route impact is

\begin{equation}
I_i(c,r,u)=\sum_k q_k(r)I_{i,k}(c,r,u)+I_{i,\mathrm{residue}}+I_{i,\mathrm{leak}},
\label{eq:impact}
\end{equation}

where $i$ indexes impact categories. The inventory includes avoided virgin production only when recovered output meets the receiving specification and actually substitutes for primary material. Downcycling is reported as a different service, not as equivalent closed-loop replacement.

\subsection{Quality-adjusted circularity}

Let $D_{k,r}$ be the fraction of recovered output that displaces a specified primary resource and $V_{k,r}$ be value or service retention relative to the functional unit. Quality-adjusted circularity is

\begin{equation}
Q(c,r)=\sum_k q_k(r)Y_{k,r}D_{k,r}V_{k,r}.
\label{eq:circularity}
\end{equation}

This is a model quantity, not a universal percentage. It is reported with the route, product grade, region, time horizon, and uncertainty interval. It avoids counting material that is collected but rejected, processed but unsold, or sold into a product that cannot deliver the original service.

\subsection{Decision and uncertainty}

Candidate $c_1$ dominates $c_2$ only when both are feasible and $c_1$ is no worse in every objective and strictly better in at least one. The nondominated set is

\begin{equation}
\mathcal{P}=\left\{c:\ g_j(c)\geq0\ \forall j,\ \nexists c'\ {\rm that\ dominates}\ c\right\},
\label{eq:pareto}
\end{equation}

where $g_j$ are service, safety, regulatory, and receiving-grade constraints. The primary report is the Pareto frontier by product and region. A policy-weighted score may be a secondary analysis after weights are frozen.

Uncertainty is decomposed across material properties, design penalties, contamination, sorting, processing, markets, regions, service failure, external costs, and interactions. Independent laboratory batches use block bootstrap. System uncertainty uses scenario sampling. Report median, interval, and high-burden tail. Missing data trigger bounded ranges or exclusion; they are not silently assigned the most favorable candidate value.

\section{Baselines, Falsification, and Safety}

\subsection{Baselines and ablations}

The comparison matrix contains the current conventional product and regional route, equal-service comparison, mass-only comparison, ideal zero-contamination recycling, realistic measured routing, design-for-recovery, quality-aware dynamic routing, reuse or refurbishment, chemical recovery where chemistry permits, and policy cases with external-cost internalization. Ablations remove identity information, disassembly design, market demand, quality grading, cycle losses, external costs, leakage, policy support, and regional infrastructure one at a time.

The framework is falsified or narrowed under the following conditions. H1 fails if technical processability does not overestimate product-grade recovery under realistic conditions. H2 fails if design-for-recovery cannot improve quality-adjusted yield without a compensating service penalty. H3 fails if quality-aware routing never improves future service or burden relative to tonnage maximization. H4 fails if reuse has no product- or region-specific cycle threshold. H5 fails if external costs never change a route ranking under credible scenarios. H6 fails if one route dominates every product class and region, which would indicate that the candidate portfolio or objective vector is incomplete.

\subsection{Controlled exclusion}

Medical, biological, battery, flame-retarded, radioactive, and chemically hazardous materials are not placed into open experimental recycling streams. They require applicable containment, occupational controls, chain of custody, and regulatory review. A controlled treatment or substitution route is a valid system result. It is not a failure to meet an unconditional recycling target.

The same principle applies to dissipative materials that have dispersed through the environment or become too dilute to recover safely. The framework records the loss and tests prevention, substitution, capture, and containment. It does not convert an unsafe recovery operation into a circularity success merely because some mass can be collected.

\section{Scale-Up, Markets, and Governance}

\subsection{From pilot yield to industrial throughput}

A process that works on clean laboratory samples can fail at industrial scale because feed composition, residence time, moisture, wear, maintenance, and product variability change together. The pilot stage therefore closes a mass and energy balance from incoming objects to saleable recovered product. For each tonne entering a facility, record collected mass, sorted mass, preprocessing loss, process yield, product-grade mass, residue, emissions, wastewater, energy, labor, and downtime. The same balance is measured for the conventional route so that capacity expansion is not credited merely because a new machine accepts more feed.

Scale-up also changes the meaning of contamination. A small amount of an incompatible additive may be harmless in one batch and unacceptable in another when it accumulates over repeated cycles. A sorting line may produce a high average purity while sending rare hazardous objects to workers or residual fractions. The design therefore reports distributional outcomes, including the lower-quality tail of batches, exposure controls, and the probability of facility shutdown. Facility throughput is optimized jointly with product grade rather than maximized independently.

\subsection{Market closure and material displacement}

Recovered material is circular only when it can enter a defined next product. A market experiment should specify the buyer, grade, test method, contract volume, acceptable contamination, and substitute material. It should compare recovered output with the marginal virgin material actually displaced, not with an abstract average. If demand disappears when virgin prices fall, the model records stranded inventory and storage or disposal burden. If recovered material is sold into a lower-grade product, the functional unit is changed and the next recovery step is tracked.

The market model varies virgin price, recovered price, transport, product demand, minimum recycled content, quality discounts, and facility utilization. It also tests coordinated contracts among producers, municipalities, processors, and buyers. A contract is successful only if it survives price and throughput stress without shifting residues or exposure to an unpriced community. Public procurement, recycled-content rules, extended producer responsibility, deposit-return schemes, disposal fees, and landfill taxes are scenario variables. Their effects are reported separately from the intrinsic process yield.

\subsection{Private and social feasibility}

Private feasibility is defined by operating cost, capital recovery, revenue, and risk. Social feasibility adds climate damage, pollution, cleanup, ecological loss, health effects, worker exposure, traffic, noise, and distributional burden. Let $C_{\mathrm{priv}}$ denote private cost and $X_j$ external burden $j$ with policy weight $\tau_j$. The scenario social cost is

\begin{equation}
C_{\mathrm{soc}}=C_{\mathrm{priv}}+\sum_j \tau_j X_j.
\label{eq:socialcost}
\end{equation}

The weights are policy assumptions and are varied transparently. The equation does not imply that all social value can be represented monetarily; non-monetized harms remain reported as constraints or separate objectives. This separation matters when an informal route has a low financial cost but high exposure, or when a formal route has a high financial cost but prevents persistent leakage.

\subsection{Governance and data quality}

Near-universal recovery requires consistent definitions across producers, waste operators, regulators, and buyers. The data layer records composition, additive hazards, repair history, contamination, route eligibility, product grade, and chain of custody. It should disclose only information needed for safe routing and measurement, with access controls for commercially sensitive data. A material passport is useful only if its entries are auditable and if the facility can act on them.

Governance experiments compare voluntary design guidance, producer responsibility, deposit-return, standardized labels, recycled-content requirements, and combinations of these instruments. Outcomes include collection, quality, cost, emissions, informal-sector displacement, worker safety, and access for low-income communities. A policy that increases nominal recycling while worsening exposure or excluding communities is not accepted as a successful circularity intervention.

\begin{table}[t]
\caption{Scale-up gates and failure interpretation}
\label{tab:gates}
\centering
\footnotesize
\setlength{\tabcolsep}{3pt}
\begin{tabular}{p{0.28\columnwidth}p{0.60\columnwidth}}
\toprule
Gate & Advance condition\\
\midrule
Process & Repeatable throughput, closed mass balance, controlled emissions, acceptable downtime and worker exposure.\\
Product & Recovered output meets a specified grade, safety requirement, and receiving-product test.\\
Market & An identified buyer and virgin displacement remain under price and demand stress.\\
System & Service, environmental, social, and cost objectives remain acceptable under regional routing uncertainty.\\
Governance & Responsibilities for design, collection, residues, data, and external burdens are enforceable.\\
\bottomrule
\end{tabular}
\end{table}

\subsection{Resilience rather than a single best case}

Material systems operate under shocks: a buyer closes, electricity becomes carbon intensive, a storm interrupts transport, contamination rises, a facility fails, or a policy changes. For scenario $s$, define route resilience as the fraction of required service retained after the shock:

\begin{equation}
R(c,r)=\min_{s\in\mathcal{S}}\frac{S(c,r,s)}{S_{\mathrm{baseline}}(r,s)}.
\label{eq:resilience}
\end{equation}

The minimum is intentionally demanding; the full distribution over scenarios is also reported. A system with a high average recovery rate but catastrophic failure under one plausible shock may require redundant facilities or a different design. Diversifying material streams can improve resilience, but it may also increase sorting complexity. The portfolio optimizer exposes this trade-off.

These scale-up gates make the central question operational. A material system can be called broadly circular only after it produces a usable grade, has a buyer, covers its relevant burdens, and remains functional when conditions depart from an ideal demonstration. If it fails one gate, the failure points to a specific intervention: better design, cleaner collection, improved sorting, more capacity, a different buyer, stronger policy, or safe substitution.

\section{Portfolio Optimization and Transfer}

\subsection{Joint selection of routes}

The decision problem is a constrained portfolio problem rather than a search for a single champion material. For each product class, select a design, an identity level, a route mix, a facility allocation, a cycle policy, and a market instrument. The optimizer maximizes quality-adjusted future service while constraining safety, greenhouse gases, energy, water, toxicity, persistence, private cost, social cost, and regional access. It may select different routes for different objects in the same nominal material family.

The model uses a two-stage procedure. First, a feasibility filter removes candidates that fail service, safety, receiving-grade, or regulatory requirements. Second, a multi-objective search identifies nondominated portfolios. A material-only model is retained as an ablation. The difference between the full and material-only Pareto sets measures the value of product and system co-design. A tonnage-maximizing policy is also retained as a baseline because high mass diversion can coexist with low quality and low future service.

Portfolio decisions are stress-tested under three classes of uncertainty. Aleatory uncertainty covers object composition, contamination, sorting outcome, breakage, and user behavior. Epistemic uncertainty covers incomplete inventories, unmeasured externalities, uncertain market displacement, and route data gaps. Structural uncertainty covers alternative assumptions about allocation, policy, leakage, and the definition of acceptable quality. The report gives both scenario-specific decisions and a stability map showing which routes remain feasible when assumptions change.

\subsection{Out-of-sample and temporal validation}

Validation must not stop at the facility that supplied the calibration data. Hold out one product variant, one facility, and one region. Freeze service thresholds, quality grades, and impact boundaries before evaluating the holdout. Transfer performance is measured by route-yield error, quality-prediction error, service-failure error, and rank stability of the Pareto portfolio. If an identity model performs well only for the products on which it was trained, the information layer requires redesign.

The system is also evaluated over time. Repeat material-flow audits after seasonal changes and after one market or policy intervention. Track whether quality, costs, and route probabilities drift as products, facility maintenance, and recovered-material demand change. An apparent improvement that disappears after a price shock or facility outage is classified as fragile. A route that remains useful across time and regions is a stronger candidate for expansion.

\subsection{Decision outputs for practitioners}

The final output is a route card for each product class and region. It states the accepted design, composition and identity requirements, contamination limit, sorting method, receiving grade, buyer, cycle or processing limit, residue route, private cost range, social burden range, and environmental conditions for advancement. A route card also states when the system should stop recovery and use controlled treatment or substitution.

This format translates the research into action without collapsing its uncertainty. Producers can redesign joins, additives, and labels. Waste operators can specify sensors, capacity, and contamination thresholds. Processors can report grades and residue. Buyers can define displacement. Regulators can compare collection, content, deposit, disposal, and producer-responsibility instruments. Communities can see worker, traffic, pollution, and distributional effects. The portfolio is therefore auditable at the point where a material claim becomes an operational decision.

\subsection{What would count as success}

Success is a measured increase in delivered future service per unit of primary resource, not merely a larger number printed on a recycling report. A portfolio passes when it satisfies functional and safety constraints, improves quality-adjusted recovery or reuse, identifies a receiving market, reduces or bounds relevant impacts, and remains viable under plausible route and market shocks. It does not need to recover every object to demonstrate that more material can circulate.

Conversely, a negative result is informative. If a product cannot pass a contamination threshold, the intervention is better source separation or redesign. If it cannot meet a receiving grade after repeated cycles, the route is downcycling or controlled treatment. If it has no buyer, the intervention is demand creation or substitution. If external burdens dominate, the process requires a different energy source, scale, containment system, or policy. These outcomes make the research useful even when the universal claim is rejected.

\section{Discussion}

The evidence supports an ambitious but conditional conclusion. Many more materials can be kept in use if products are designed for disassembly, material identity is available, contamination is controlled, facilities are accessible, output quality is specified, and buyers exist. The combination of robotics, spectroscopy, selective chemistry, digital product information, producer responsibility, and stable recovered-content demand can expand circularity well beyond current practice.

However, the phrase "every material" is not a useful universal engineering specification. Recovery becomes physically and economically difficult when materials are mixed, diluted, contaminated, degraded, chemically transformed, or dispersed. A high collection rate can conceal low quality. A high process yield can conceal solvent, energy, toxicity, or residue burdens. A profitable route can externalize health or cleanup costs. A reuse claim can conceal washing and replacement. These are measurable trade-offs, not reasons to abandon innovation.

The most effective target is therefore a portfolio of material-product systems. Clean and durable objects should first be considered for reuse, repair, or refurbishment. Narrow and compatible streams should be directed to mechanical recycling. Chemically selective routes should be used where recovered feedstock actually displaces virgin material. Contaminated or hazardous streams should receive controlled treatment or substitution. Design, identity, route, market, and policy should be optimized together.

This conclusion is stronger than a generic call for more recycling. It identifies a practical decision rule: a route advances only when it passes service and safety gates, produces an identified grade and buyer, and retains quality-adjusted benefit under plausible regional scenarios. A candidate can fail and still produce actionable knowledge by revealing the required infrastructure, cycle count, energy threshold, contamination limit, or redesign.

\section{Conclusion}

The prospect of recycling and reusing essentially every material should be separated into two claims. Near-universal circularity for a designed portfolio of products is a credible research and engineering objective. Literal repeated recovery of every object is not credible as a blanket promise because some streams are hazardous, contaminated, dissipative, mixed, or too costly to recover safely.

UNIVERSAL-CIRCULAR-RECOVERY provides a testable route to the first objective. It compares product architecture, material chemistry, information, sorting, recovery, markets, and policy per delivered service. It measures quality retention, repeated service, contamination, residues, leakage, energy, water, environmental burden, private cost, social cost, and resilience. It reports Pareto frontiers and conditional routes rather than a single recycling percentage.

The direct answer is yes for many carefully designed material-product systems, but not because a universal machine will process every object. It will require reduction where possible, reuse and repair where durable, high-quality recycling where streams are compatible, selective chemical recovery where it truly displaces virgin feedstock, controlled treatment for hazardous materials, and product and policy design that make the loop economically and socially real.

\begin{thebibliography}{00}
\bibitem{epa} U.S. Environmental Protection Agency, \textquotedblleft Recycling Basics and Benefits,\textquotedblright\ updated Jun. 22, 2026. [Online]. Available: \url{https://www.epa.gov/recycle/recycling-basics-and-benefits}
\bibitem{geyer} R. Geyer, J. R. Jambeck, and K. L. Law, \textquotedblleft Production, use, and fate of all plastics ever made,\textquotedblright\ \textit{Science Advances}, vol. 3, no. 7, art. no. e1700782, 2017, doi: 10.1126/sciadv.1700782.
\bibitem{hopewell} J. Hopewell, R. Dvorak, and E. Kosior, \textquotedblleft Plastics recycling: challenges and opportunities,\textquotedblright\ \textit{Philosophical Transactions of the Royal Society B}, vol. 364, no. 1526, pp. 2115--2126, 2009, doi: 10.1098/rstb.2008.0311.
\bibitem{iso18604} International Organization for Standardization, \textit{ISO 18604:2013 Packaging and the Environment--Material Recycling}, Geneva, Switzerland, 2013.
\bibitem{iso15270} International Organization for Standardization, \textit{ISO 15270:2008 Plastics--Guidelines for the Recovery and Recycling of Plastics Waste}, Geneva, Switzerland, 2008.
\bibitem{oecd} Organisation for Economic Co-operation and Development, \textit{Global Plastics Outlook: Economic Drivers, Environmental Impacts and Policy Options}, OECD Publishing, Paris, France, 2022, doi: 10.1787/de747aef-en.
\bibitem{allwood} J. M. Allwood, \textquotedblleft Squaring the circular economy: the role of recycling within a hierarchy of material management strategies,\textquotedblright\ in \textit{Handbook of Recycling}, E. Worrell and M. A. Reuter, Eds. Amsterdam, The Netherlands: Elsevier, 2014, pp. 445--477.
\end{thebibliography}

\end{document}
